\subsubsection{Smaller jumps and the two estimates at ordinary inert primes}
\label{sec:inert-row-table}
\hypertarget{source-inert-row-table}{}

When smaller jumps occur, the uniform restriction used for the larger jumps does not cover the whole sequence. Decomposing it into short arcs gives the missing estimate; it can then be combined with the larger-jump case at the same inert place.

Here \(\ell\ge11\) and \(n=n_\lambda\) is the local unramified degree. In the
case with some smaller jumps, let \(s_i\) be the sizes in the local word.
Use \(b_0,w,S,\eta\) for its threshold for the larger jumps, bound for the smaller jumps,
spread and central width. Each
\[
 s_i>\eta.
\]
A size in the larger range is
in \([b_0,b_0+\eta]\), and a size in the smaller range is at most \(w\).
Use the signed jump intervals of Section~\ref{subsec:inert-rows}.

At every letter at most \(m\)
target jumps, counted with multiplicity, lie in \([-\eta,\eta]\).
On the marked row the signed windows come from the source lattice
in the sum of twisted duals. The bound on the central multiplicity
comes from the twisted dual target at each labeled embedding,
applied with the same enumeration of the embeddings as the word.
Both statements therefore apply to a size in the smaller range.
\label{cor:inert-row-table-1}
There is a position with smaller jumps, the signs alternate cyclically,
and there are at most two consecutive positions with larger jumps.
The numerical hypotheses are
\[
 20w,20S,20\eta\le b_0.
\]

Put \(h=b_0+\eta+S\) if there is a position with larger jumps and \(h=w+S\)
otherwise, and put
\[
 \gamma_h=\log_\ell(2h),\qquad
 S_g=h/\ell^2+(w+S+\eta)/\ell+1+\gamma_h .
\]
Let \(B\) be the indicator that some position with smaller jumps \(i\)
has a neighbor with larger jumps and satisfies
\[
 \abs{\ell s_i-b_0}\le 320
              (\ell\eta+\ell S+b_0/\ell+s_i+\ell\gamma_h).
\]
The common place indicator in \(F_{\ell,\lambda}\) is the
existence of such a word on the marked data, over all good
characters and compatible local coefficient data. Therefore every
local \(B\) is at most this place indicator.

Cut before each position with smaller jumps, numbering the cut by
its preceding letter. Begin the enumeration at one such position, so
the construction includes the arc that crosses the end of the cycle.
Let \(\mathcal C\) be the set of cuts. The arc ending at \(c\)
consists of one smaller letter followed by \(k_c\) larger letters.
Consequently
\[
 a_c=1+k_c,\qquad 0\le k_c\le2,\qquad
 \sum_{c\in\mathcal C}a_c=n .
\]

Suppose first that a larger letter occurs and \(40S_g>b_0/4\).
The two jump bounds of the coarse lift are at most \(h\); its
maximum width is at most \(h+\eta\). The coarse exponent calculation
in the proof of Proposition~\ref{prop:ordinary-inert-place-entry}
therefore bounds the exponent \(e'\) of a full finite lift by
\(7.1 h/\ell+7.1+2\gamma_h\). Since \(h\le1.1b_0\),
\(b_0<160S_g\), and \(\ell\ge11\),
\[
 7.1 h/\ell <113.6 S_g,\qquad
 7.1+2\gamma_h\le7.1S_g,\qquad e'<160 S_g .
\]

In the remaining range, use Lemmas~\ref{lem:two-letter-wach-lift},
\ref{lem:three-letter-negative-lift} and~\ref{lem:three-letter-positive-lift}
on arcs of length two and three. To compare the arc lift with the cut
decomposition we need the exponents \(R+2\sigma\) and \(R+3\sigma\).
Put \(g=w+S+\eta\), \(q_0=\ell-1\), and
\(M=\lfloor2h/q_0\rfloor+1\). There is one smaller letter
immediately after the cut. The corresponding correction satisfies
\[
 \begin{aligned}
 \sigma
 &=\left\lceil1+\frac{\ell(2(w+S)+M)}{q_0^2}\right\rceil\\
 &\le 2+\frac{\ell}{q_0^2}
       +\frac{2\ell g}{q_0^2}+\frac{2\ell h}{q_0^3}
 \le 3\left(1+\frac{g}{\ell}+\frac{h}{\ell^2}\right).
 \end{aligned}
\]
Indeed \(\ell/q_0\le11/10\) bounds the two variable coefficients
after normalization by \(2(11/10)^2\) and \(2(11/10)^3\).
For the recurrence of Section~\ref{sec:inert-long-orbit}, take
symmetric widths at most \(2h\) throughout and at most \(2g\)
at the first letter after the cut. With
\(c_H=\lceil1+\log_\ell(1+\max_i H_i)\rceil\) one has
\[
 R\le c_H + \frac{\ell}{q_0(q_0-1)}
                    (2g+c_H+3+2h/q_0).
\]

Here \(c_H\le 2+\tfrac15+\gamma_h\) by
\(1+2h\le3h\) and \((3/2)^5<11\).
The coefficients of \(g/\ell\) and \(h/\ell^2\) in the bound
are at most \(2\cdot11^2/(10\cdot9)<3\) and
\(2\cdot11^3/(10^2\cdot9)<3\), respectively. Since
\(\ell/(q_0(q_0-1))\le11/90\), the remaining terms
are at most \(3(1+\gamma_h)\). Combining these estimates gives
\begin{equation}\label{eq:inert-short-cut-budget}
 R\le3S_g,\qquad \sigma\le3S_g,\qquad
 R+2\sigma\le9S_g,\quad R+3\sigma\le12S_g .
\end{equation}
With the cut costs of \eqref{eq:inert-short-cut-budget}, the arc lift at \(c\) itself has exponent at most
\[
 \sigma+\kappa_c+\lfloor 40S_g+\eta B\rfloor ,
 \qquad
 \kappa_c\le 13S+3(1+\gamma_h).
\]
It lands in the finite image modulo a span of at most
\(m r_c(a_c+1)\) target finite classes, with \(r_c\le a_c\).
We derive the coordinate lift and the fixed exceptional span below.

\paragraph{Coordinates on a short arc.}\label{par:short-wach-frame}
\begin{proposition}[Common coordinates for short-arc comparisons]\label{prop:short-arc-coordinates}
For a short arc, use the same integral vertical and
horizontal frames and the corresponding cyclotomic levels.
For an arc ending at \(c\) of length \(a=a_c\), the offset \(0\le j<a\)
has size \(u_j=s_{c-j}\) and cyclotomic factor \(P_j=\phi^j(X+\ell)\). The horizontal and vertical comparisons use the same positive pole and negative zero factors
\[
 D=\prod_{j\ {\rm positive}}P_j^{u_j},\qquad V=\prod_{j\ {\rm negative}}P_j^{u_j+S}.
\]
The two ordinary lifting rules use these coordinates. In the single-letter edge cases the
positive rule applies to every input, and the negative rule
uses the empty pole set and the initial constant lift.
\end{proposition}

\begin{proof}
\proofstep{Fix the same cut coordinate and ordered cyclotomic levels for both rules.}
Write \(a=a_c,\ c'=c-a,\ G=w+S,\ u_j=s_{c-j}\), and
\(P_j=\phi^j(q)\), so \(P_0=q=X+\ell\).
Let \(\mathcal R=O[[X]]=O[[q]]\). All ideals of fractions below
are in the component at \(c\).
Call an offset positive or negative according to its large Hodge
jumps. This sign is used even at the first letter of small size.
The full denominator for transport over \(a\) steps and the two
smaller products determined by the jumps are
\[
 \Lambda=P_{a-1}^{G}\prod_{j<a-1}P_j^h,\qquad
 D=\prod_{j\ {\rm positive}}P_j^{u_j},\qquad
 V=\prod_{j\ {\rm negative}}P_j^{u_j+S}.
\]
Choose a basis of the tame-invariant Wach lattice at \(c'\).
The unique formal \(\gamma\)-fixed basis with the same constants is
its horizontal basis.

Write \(U(X)=I+\sum_{r>0}U_rX^r\) for the change
to the horizontal basis, and let \(F_r(U_1,\ldots,U_{r-1})\)
be the degree-\(r\) residual computed with \(U_r=0\).
The recursion is
\[
       (\omega_X^r-1)U_r=-F_r(U_1,\ldots,U_{r-1}).
\]
Its new denominator has order \(1+v_\ell(r)\) at step \(r\).
Thus through that step the denominators divide
\(\ell^{r+v_\ell(r!)}\). This also holds for the inverse change
of basis, by the factorial inequality for products of coefficients.
Use the basis truncated before degree \(M\). These are coordinates
on the given lattice map: a map of Wach lattices acts on horizontal
columns by its matrix of constants at \(c'\), since both sides solve
the same equation with the same constant. Frobenius transports these
constants to the crystalline space at \(c\).

Here are details on the precision of these coordinates. Put
\[
 P(T)=\langle\Phi^aN_{c'}(T)\rangle_{\mathcal R}.
\]

Represent \(v\in\Lambda^{-1}P(T)\) modulo \(X\Lambda P(T)\)
in the transported horizontal basis, allowing powers of \(\ell\)
to clear coefficients in this comparison. Equivalently multiply
by \(\Lambda\) and compare modulo \(X\Lambda^2 P(T)\).
The tails beyond the truncation have a factor
\(\phi^a(X)^M\), divisible by \(X\Lambda^2\) since
\((\ell-1)M>2h\). The finite coefficients, and the inverse matrix,
can be transported integrally at cost \(\sigma\). Indeed on
the first step one reduces \(\ell^\sigma\phi(X)^r/\ell^{r+v_\ell(r!)}\)
modulo \(X^Mq^{2G}\).

The term \(X^r q^{(\ell-1)r}\) pays the
denominator after this precision, by the binomial division of
Paragraph~\ref{par:wach-cut-correction} with one forward step and
precision \(2G\). More explicitly, for \(k=r+v_\ell(r!)-\sigma>0\)
one needs
\[
 (\ell-1)r\ge 2G+k+M-r-1.
\]
Then modulo \(\ell^k\)
the binomial expansion of \((X+\ell)^{(\ell-1)r-2G}\) is divisible
by \(X^{M-r}\). If \(k>0\), the coefficient bound gives
\(r>(\ell-1)\sigma/\ell\). Together with the definition of
\(\sigma\) this yields the preceding power inequality.

In
the binomial expansion the terms of \(\ell\)-degree \(\ge k\)
vanish modulo \(\ell^k\); every other term has \(X\)-degree
at least \((\ell-1)r-2G-k+1\ge M-r\). Thus the termwise
division is legitimate at the first precision. The next
\(a-1\) Frobenius maps give
\[
             \phi^{a-1}\bigl((X^Mq^{2G})\bigr)
                                     \subseteq (X\Lambda^2)
\]
by the same level bounds defining \(\Lambda\). Thus the transported matrices are represented by
integral lifts of \(\ell^\sigma\phi^a(U)\) and
\(\ell^\sigma\phi^a(U^{-1})\) modulo \(X\Lambda^2\)
in the Frobenius basis. If \(A,B\) denote these two integral lifts in the displayed order, their common precision gives
\[
 AB\equiv BA\equiv\ell^{2\sigma}I\pmod{X\Lambda^2}.
\]


For clarity, the coefficient test in these coordinates is as
follows. If a constant direction \(d\) at the end of the arc
belongs to \(\mathcal F_j^{-\alpha_j}\) for each transported
filtration, then \(d\) with coefficient in
\((\prod P_j^{\alpha_j})\mathcal R\) passes the Wach test at
each level. Sums of such terms after scalar extension give the
rational test; the bounded integral numerators are required as
well.

To see the rational assertion, separate the levels by the
Chinese remainder theorem for their distinguished polynomials
and take a basis of the filtration at each level. For a basis vector \(d\) of jump \(u\), the corresponding generator and its Laurent order are
\[
 t^{-u}d,\qquad \operatorname{ord}_t(t^{-u})=-u,
\]
as in
Paragraph~\ref{par:wach-level-models}. Multiplying the other
level factors to sufficiently high orders realizes each of these
finite jets without imposing further conditions there.
Divisibility of bounded numerators by a product of these
polynomials descends to \(\mathcal R\) by monic division.

Let \(u\) denote a coordinate vector satisfying all the
local tests and let \(n\) be a vector in \(N_c\).
The forward comparison and its inverse give, in their
respective quotients by the transported precision,
\[
   f_{\rm hor}(u)=[\ell^\sigma u]_{\rm Wach},
   \qquad g_{\rm hor}(n)=[\ell^\sigma n]_{\rm coordinates},
   \qquad g_{\rm hor}f_{\rm hor}=\ell^{2\sigma}.
\]
Thus an admissible coordinate vector produces a Wach class,
and the restricted inverse sends a Wach vector back into the
tested coordinate module. Indeed the transported forward
and inverse matrices above give integral lifts with the bounded
denominator, and changes by \(X\Lambda P(T)\) lie in \(XN_c\)
by the height bounds.

For an admissible vector one tests
\((\gamma-1)/X\) in the same manner. This gives an integral
vector with \((\gamma-1)v\in XN_c\), since its horizontal
columns are fixed to the indicated precision. A product of the
\(P_j\) clears its denominators in \(N_c\). These statements
apply to subarcs beginning at \(c'\) with their own product
\(\Lambda\). \label{cor:inert-short-frame-1}

We will lift a coordinate \(f\) on a target horizontal column
\(d\) which belongs to \(\mathcal F_j^{u_j}\) at every positive
level. Its scalar hypotheses are
\begin{equation}\label{eq:short-scalar}
 Df\in\mathcal R,\qquad
 D(\ell^{k_j}f-x_j)\in P_j^{u_j+S}\mathcal R
 \quad(j\ {\rm negative})
\end{equation}
for some \(k_j\ge0,\ x_j\in O\). A lift at cost \(e\) represents
\(\ell^e f d\) by the image of an admissible source vector up to
\(XN_c(T_2)\) and the coordinate precision. There are two ways
to construct it.
\begin{enumerate}
\item If \(Df\in V\mathcal R\), then \(fd\) passes the target
Wach test. The poles at the positive levels are allowed by the
filtration of \(d\); the zeros at the negative levels accommodate
the lowest possible jump. The coordinate comparison and scalar image bound at \(c\) use successive multipliers
\[
 \ell^\sigma\ell^R=\ell^{R+\sigma}.
\]
Tame averaging preserves the lifted image, so no extra exponent enters this case.
\item Suppose the poles are confined to a subset of the positive
levels and an integral source direction maps to \(\ell^t d\) and
lies in \(\mathcal F_j^{u_j}\) there. At each omitted positive
level also require that \(f\) is constant modulo \(P_j^\eta\)
after clearing scalars and the retained denominator. Then the
coefficient \((\gamma f-f)/X\) times this source direction passes
the level tests. For each level factor, write \(\gamma P_j/P_j=1+Xh_j\) with \(h_j\in\mathcal R\). Then
\[
 (\gamma P_j-P_j)/X=h_jP_j.
\]

For every \(n\in\mathbb Z\) the ratio \((1+Xh_j)^n\) belongs to \(1+X\mathcal R\). Thus
\[
 (\gamma P_j^n-P_j^n)/X\in P_j^n\mathcal R ,
\]
including for negative powers. The constant disappears
at each omitted positive or negative level. Equation~\eqref{eq:short-scalar}
therefore supplies the required zeros at negative levels, and
\(\eta\) supplies them at omitted positive levels. The retained
positive poles pass by the source filtration. The admissible-vector
construction gives the lift at \(t+\sigma\). For an integral constant, transport the constant term of the initial target basis lift at \(c'\). Its coefficient differences vanish, so its cost is at most \(R+\sigma\).
\end{enumerate}

In the second rule the congruences use the denominator for the
retained positive levels. Costs of two summands may be padded
to their maximum, by linearity of the source map and of the
coordinate comparison. For a single positive letter the first
rule applies to every \(f\). For a single negative letter use
the second with empty pole set and the initial constant lift.
\end{proof}


\begin{auxlemma}[Separating a scalar by Weierstrass division]\label{lem:short-arc-content}
Let \(O\) be a coefficient DVR with fraction field \(K\),
\(v_p(p)=1\), and let \(J\subset\mathcal R=O[[q]]\) be generated by
a distinguished polynomial. Take \(D\in\mathcal R\) and
\(f\in\operatorname{Frac}(\mathcal R)\). Assuming that the residue
\(\bar D\in\mathcal R/J\) is nonzero, denote by \(c_J(D)\) the minimum of the
valuations of its Weierstrass remainder coefficients. Let \(k,d\)
be nonnegative integers with \(c_J(D)\le d\). If
\(Df\in\mathcal R\), \(x\in K\), and \(p^k Df-Dx\in J\), then there is an integral constant satisfying
\[
 \exists x'\in O:\qquad D(p^d f-x')\in J.
\]
\end{auxlemma}
\begin{proof}
Write \(H=\mathcal R/J\), a free \(O\)-module with the monomials
below \(\deg J\) as basis, where \(\deg J\) is the degree of the
distinguished generator. In \(H\otimes_O K\) the given equality
implies
\[
 x\bar D=p^k\overline{Df}\in p^k H,\qquad
 v_p(x)+c_J(D)\ge k .
\]
Thus \(x'=p^{d-k}x\) is integral, including when \(k>d\).
Both \(p^d Df\) and \(x'D\) lie in \(\mathcal R\), and their
images agree after tensoring \(H\) with \(K\). Freeness makes
\(H\to H\otimes_O K\) injective, which proves the assertion.
\end{proof}

\begin{lemma}[Two-letter lifting on a short arc]\label{lem:two-letter-wach-lift}
In the models of Proposition~\ref{prop:short-arc-coordinates},
write \(e_2\) for the exponent of a two-letter correction. One can achieve
\[
 e_2\le15L+1+\eta B,
\]
where \(L\) is the local size from the short-arc comparison and \(B\) its
exceptional indicator.
\end{lemma}

\begin{proof}
\label{par:short-two-levels}
Use \(u=s_{c-1}\) for the small size and \(v=s_c\) for the large
one. Put \(p=\ell,\ L=S_g\) within this calculation. Denote by
\(A_i\) the ceiling of the least ordinary Wach image exponent in
\(\pi\) at \(i\), divided by \(p-1\) before taking the ceiling.
Then \(X^{A_i}\) kills the tame-invariant cokernel. In signed
boxes of widths \(s_i+S+\eta\), the recurrence of
Paragraph~\ref{par:wach-cut-colength} gives
\begin{equation}\label{eq:short-A}
 A_i\le(s_{i+1}+S+\eta)/p+\tfrac52 L.
\end{equation}
Here one can replace \(5/2\) by \(2\) if \(i+2\) is small.

\proofstep{Compute the denominators of the two neighboring level evaluations.}
To check the constants, write \(d=p-1,\ H=h+\eta\),
\(W_j=s_j+S+\eta\) and \(c_H=\lceil1+\log_p(1+H)\rceil\).
The ordinary image exponents satisfy
\(s_i^{\rm im}\le W_{i+1}+c_H+3+s_{i+1}^{\rm im}/d\),
and are at most \(H\) by the signed twist. Thus
\[
 A_i-W_{i+1}/p\le 1+(c_H+3)/d+H/(pd)+H/d^2.
\]
The general bound uses \(H=h+\eta\), \(H+1\le2h\), and
\(p/d\le1.1\).

If \(i+2\) is small, iterating the recurrence
once replaces its last summand by
\[
 \frac{H}{d^2}\quad\rightsquigarrow\quad
 \frac{w+S+\eta+c_H+3}{d^2}+\frac{H}{d^3}.
\]
The respective coefficients of \(h/p^2\) are bounded by
\[
 \frac{p}{d}+\frac{p^2}{d^2}\le2.31,\qquad
 \frac{p}{d}+\frac{p^2}{d^3}\le1.24 .
\]
The \(\eta/p^2\) terms are absorbed in \((w+S+\eta)/p\).
Together with \(c_H\le2+\log_p(2h)\), these give the factors
\(5/2\) and \(2\) in the two versions of \eqref{eq:short-A}.
\label{cor:inert-short-two-1}

Apply Lemma~\ref{lem:short-arc-content} with
\(P=P_1=p\varepsilon+Uq^p\), where \(\varepsilon,U\in\mathcal R^\times\).
We estimate the contents for the two possible denominators.
For \(D=q^v,J=(P^y)\), put \(t=\lfloor v/p\rfloor+1\).
Multiply by \(q^{pt-v}\) and expand
\[
 (q^p)^t=U^{-t}(P-p\varepsilon)^t.
\]
Modulo \(P^y\) the
content of this power is the minimum over \(j<y,\ j\le t\)
of \(t-j+v_p\binom tj\). After division by the lowest
uniformizer power, the least \(j\) attaining it has strictly
smaller order in the residue parameter than the other surviving
terms and than \(P^y\). This proves the minimum formula even for
nonconstant units.

Since multiplication by \(q^{pt-v}\) can only
increase content in the first quotient, any admissible index gives
an upper bound for \(c_{(P^y)}(q^v)\).
For \(D=P^u,J=(q^C)\), use the binomial expansion of
\((p\varepsilon+Uq^p)^u\) modulo \(q^C\). The two bounds are
\begin{equation}\label{eq:two-letter-wach-contents}
 \begin{aligned}
 c_{(P^y)}(q^v)&\le t-i_1+v_p\binom{t}{i_1},
       & i_1&=\min(t,y-1),\\
 c_{(q^C)}(P^u)&\le u-i_2+v_p\binom{u}{i_2},
       & i_2&=\min(u,\lfloor(C-1)/p\rfloor).
 \end{aligned}
\end{equation}
For \(pu<C\) the second line has \(i_2=u\) and the content is zero.
For the binomial terms in \eqref{eq:two-letter-wach-contents}, the carry formula gives \(v_p\binom ti\le\log_p t\).
In particular the two bounds are at most
\[
 \max(0,v/p-y+2+\log_p v)+1,\qquad
 \lceil\max(0,u-C/p+2+\log_p u)\rceil ,
\]
respectively. In the first bound the case \(i=t\) contributes
zero. All contents are normalized by \(v_p(p)=1\).

Suppose first that the large terminal letter is positive. If
\(p(u+S+\eta)+\eta<v\), construct a filtered source constant.
On the target take the coordinate \(P^{u+S}/q^v\) in direction
\(d\), realize it at cost \(\sigma\) and lift at the cut at
cost \(R\). Express the lift in source horizontal coordinates
at the cost of a further \(\sigma\). Multiply these coordinates
by \(q^{v+S}P^\eta\), obtaining integral series. In fact every
nonzero filtered summand has exponents at least \(-(v+S)\)
and \(-\eta\) at the two levels. Divisibility after rational
expansion descends by saturation. In the coordinate \(q\) take
the coefficient of degree
\[
 j=S+p(u+S+\eta).
\]

A filtered summand with level-zero exponent \(\alpha\) can
contribute only if
\[
 v+S+\alpha\le j.
\]
This forces
\(-\alpha\ge\eta+1\); the gap puts its direction in
\(\mathcal F_0^v\). The constant matrix of the map sends this
coefficient to the same coefficient of
\(p^{R+2\sigma}q^S P^{u+S+\eta}d\). Indeed the coordinate error
after clearing these denominators is divisible by \(q^h\) and
\(j<h\). The coefficient of degree \(p(u+S+\eta)\) of
\(P^{u+S+\eta}\) is a unit, by reduction modulo the maximal
ideal. Dividing by this unit supplies the filtered preimage of \(p^{R+2\sigma}d\). The additional coordinate comparison in the second rule raises the exponent to
\[
 (R+2\sigma)+\sigma=R+3\sigma.
\]


In the complementary range take \(D=q^v\), \(y=u+S\). The
first content bound separates the constant at the negative level
from the remainder. Each is lifted at cost \(R+\sigma\) by the
two rules. Writing the extra cost as \(c_{\rm cont}\), the two
ranges give
\[
 c_{\rm cont}\le
 \begin{cases}
    3+\gamma_h,& v<py,\\
    \eta+\eta/p+3+\gamma_h,& py\le v\le p(y+\eta)+\eta .
 \end{cases}
\]
In the second range, combining the two bounds on \(v\) with
\(b_0\le v\le b_0+\eta\) gives
\[
 b_0/p-S-\eta-\eta/p\le u\le(b_0+\eta)/p-S,
\]
the prescribed band.

Now suppose the first small letter is positive. Put
\(A=A_{c-1}\). When \(u\ge A+\eta+1\), take the partial
arc of length one at this letter and the coordinate \(q^{-u}d\)
there, with \(d\) read before the last Frobenius transport.
Realize it at cost \(\sigma\), lift \(X^A\) times this integral
vector, and read the source coordinates at cost \(\sigma\).
Multiply by \(q^G\) and take the coefficient in degree \(G-u+A\).

The degree \(G-u+A\) is below \(G-\eta\), so the degree and gap test places the coefficient in \(\mathcal F^u\). The target coefficient in the chosen degree is one, as
\[
 [q^{A+G-u}](q-p)^A q^{G-u}=1.
\]
The precision includes \(q^{2G}\) after denominators are cleared, so the image of the coefficient is exactly \(p^{2\sigma}d\). Tame averaging costs a unit.
Transport forward and apply the second rule at pole set
\(\{1\}\). The cost is \(3\sigma\).

When \(u\le A+\eta\), apply the second content bound with
\(D=P^u,\ C=v+S\), and lift the separated remainder and
constant at cost \(R+\sigma\). There is no extra cost for
\(pu<C\). Otherwise \eqref{eq:short-A} with its improved
constant gives
\[
 u-C/p\le\eta+\eta/p+2L.
\]

If the band is absent, \(pu\ge C\ge b_0\) and
\(pu-b_0>2p\eta\). But
\[
 u-b_0/p\le \eta+(2\eta+S)/p+2L,
\]
so \(\eta<4L\) using
\((2\eta+S)/p\le L\) (since \(w>\eta\)).
Thus the two-letter exponent in all cases is at most
\(15L+1+\eta B\). \label{cor:inert-short-two-2}
\end{proof}


\begin{lemma}[Three-letter lifting with negative outer signs]\label{lem:three-letter-negative-lift}
With \(L=S_g\) and \(B\) the short-band indicator of Section~\ref{sec:inert-row-table},
for the chronological signs \(-,+,-\), one can choose a three-letter correction in the short-arc frames with exponent \(e_3\) satisfying
\[
 e_3\le22L+1+\eta B
\]
in the given DVR normalization.
\end{lemma}

\begin{proof}
\label{par:short-three-levels}
Put
\[
 u=s_{c-2},\ v=s_{c-1},\ z=s_c,
\]
so \(u\le w\)
and \(v,z\in[b_0,b_0+\eta]\). Here \(40L\le b_0/4\).
At a partial arc use level numbers relative to its own endpoint.

\proofstep{Start with the middle positive letter and retain the two negative neighbors.}
First consider signs \(-,+,-\) in chronological order.
If \(p(u+S+\eta)\le v\), work at the middle letter and
write \(A=A_{c-1}\), \(y=u+S+\eta\).
As in the two-letter coefficient extraction use the coordinate
\(P_1^{u+S}/q^v\) in direction \(d\). At this intermediate
point the scalar image bound at a cut is unavailable. Instead
lift \(X^A\) times the realized vector. After the two horizontal
comparisons and multiplication by \(q^{v+S}P_1^\eta\), the
scalar on the target, in the variable \(q\), is
\[
 F=p^{2\sigma}(q-p)^A q^S P_1^y .
\]
Below degree \(v+S-\eta\) every source coefficient lies in
\(\mathcal F^v\) by the gap. To recover an integral direction
it suffices that some such coefficient of \(F\) divides a small
power \(p^t\) in \(O\).

Take
\[
 t=1+\left\lfloor\tfrac{11}{10}
 (2\sigma+(\eta+A)/p+2+\gamma_h)\right\rfloor .
\]
Then \(2\sigma<t\le10L\), \(A+t+\eta\le v\).
For the first bound use \(\sigma\le3L\) and
\(A\le (h+\eta)/(p-1)+1\), giving
\[
 (\eta+A)/p\le 1.1(h+\eta)/p^2+(\eta+1)/p.
\]

Substitution in the floor bound gives \(t\le10L\).
For the second, \(A\le (b_0+2\eta+S)/(p-1)+1\),
\(10L\le b_0/16\) and \(20\eta,20S\le b_0\).
If all the indicated coefficients were divisible by \(p^t\),
put \(k=t-2\sigma\). Modulo \(p^k\) one has
\(q^{A+k-1}\in((q-p)^A)\) by the binomial expansion.
Cancel \(p^{2\sigma}\), remove the initial power of \(q\),
and use this divisibility. The result would be
\[
 P_1^y\in(p^k,q^{v-\eta-(A+t-1)}).
\]
The second content calculation, before rounding, gives
\[
 k\le\max\{0,\ y-(v-\eta-A-t+1)/p+2+\log_p y\}.
\]
Indeed using the index \(\min(y,\lfloor(C-1)/p\rfloor)\)
costs at most \(y-C/p+1+1/p+\log_p y\) when positive.

Since \(py\le v\) and \(p/(p-1)\le11/10\), the definition
of \(t\) contradicts the last display. Choose a coefficient
of valuation less than \(t\), and multiply its source vector by
the integral quotient of \(p^t\) by that coefficient.
Moving one step forward gives the positive-level preimage.
The sole pole of \(f\) is at this level, so the second lifting
rule gives cost
\[
 t+\sigma\le13L.
\]


Suppose next that \(v<p(u+S+\eta)\). The down-level constants
in \eqref{eq:short-scalar} can be different. Eliminate the one
at level two first. The denominator there is \(P_1^v\), the
modulus is \(P_2^{u+S}\). Put \(y=u+S\). Since
\(P_2=p\varepsilon_1+U_1P_1^p\) with unit coefficients, the
first content calculation (using \(P_1,P_2\) in place of
\(q,P_1\)) gives a splitting
\[
 t_2f=x_2+f_2,\quad D f_2\in P_2^{y}\mathcal R,\quad
 D(p^{k_0}f_2-x_0)\in q^{z+S}\mathcal R,\qquad D=P_1^v .
\]

If \(v<py\), the content of \(P_1^v\) modulo \(P_2^y\)
is zero by reduction and comparison of orders. Otherwise use
the ceiling multiple \(p\lceil v/p\rceil\) in its place.
For \(py\le v<p(y+\eta)\), the expansion in \(p,P_1^p\) of \(P_2^y\) costs
at most \(\eta+2+\gamma_h\), and the small letter belongs to
the band. Thus \(v_p(t_2)\le(\eta+2+\gamma_h)B\).

Write \(C=z+S\). For \(x_0\ne0\), the second congruence gives
\begin{equation}\label{eq:short-last-constant}
 v_p(x_0)-k_0\ge y-v+C/p-C/p^2-2-\gamma_h .
\end{equation}
Indeed \(P_2^y\) has content at least \(y-C/p^2\) modulo
\(q^C\). Expand first in \(p,P_1^p\), then in \(p,q^p\):
a term with less scalar valuation has \(q\)-degree at least
\(C\). The content of \(P_1^v\) modulo \(q^C\) is at most
\(v-C/p+2+\gamma_h\) by the second calculation, since
\(C\le pv\). Comparing contents of
\[
 x_0P_1^v\equiv p^{k_0}P_2^y U\bmod q^C
\]
proves the claim.

Here is how to use this scalar inequality with the middle
filtered lattice. Transport the truncated horizontal column to
the middle point, paying \(\sigma\) at the first small letter,
and denote its integral representative by \(\lambda\). It
represents \(p^\sigma d\) there.

Set \(s=P_1^y/q^v\) at this
point. The vector \(s\lambda\) lies in the target Wach lattice.
The level test uses \(d\in\mathcal F^v\) at the middle and
the zero of order \(y\) at the initial negative letter.
In addition
\[
 \gamma(s)(\gamma-1)\lambda\in Xq^{a_*}N_{c-1}(T_2)
\]
where \(a_*=A_{c-1}+r_{c-1}\) and \(r_i\) is the least
scalar Wach image exponent at \(i\). Indeed the one-step
horizontal correction has defect in \(X^M q^{2G}\) in the
first Frobenius span. The next transport gives factors
\(X^M q^{(p-1)M}P_1^{2G}\); after multiplication by
\(\gamma(s)\) these pay \(Xq^{a_*}\) and the signed height losses
\(q^\eta P_1^{u+S}\), since \((p-1)M\ge v+a_*+\eta\).

The powers required in the ordinary parameter and the scalar
give the \(q^{a_*}\) image by the binomial theorem.
The two-filter argument of Paragraph~\ref{par:thick-filtered-comparison}
therefore lifts \(p^E s\lambda\) integrally, where
\[
 E=\left\lfloor\frac{3a_*+S+2\eta-2}{p-1}
               +\log_p(v+\eta-1)\right\rfloor_{\!+}\le8L .
\]
For this application the lower filter has indices \([v-a_*+1,v]\)
and the other has indices in \([v+1,v+S+a_*]\cup[1-\eta,\eta+a_*]\).
The factors adjacent to \(j\) in the first two intervals
give factorials of combined length \(2a_*+S-1\).

In the remaining
interval of length \(2\eta+a_*\), isolate a term of largest
valuation; the others are bounded by factorials of combined
length \(2\eta+a_*-1\). The isolated valuation is at most
\(\log_p(v+\eta-1)\). If \(a_*=0\) the map is already onto
at this letter. The global ordinary width bound \(H\le h+\eta\)
and degree lowering give
\[
 a_*\le H(2p-1)/(p-1)^2+\lceil1+\log_p(1+H)\rceil+1.
\]

These imply \(2a_*+\eta+1\le v\) under \(40L\le b_0/4\)
and give the displayed estimate for \(E\), by \(p\ge11\) and
\(H\le b_0+2\eta+S\). In particular
\[
 v+a_*+\eta\le2v\le2h<(p-1)M,
\]
as used in the defect estimate. The separation gives the
disjointness required by interpolation. \label{cor:inert-short-three-1}

Let \(n_1\) be this lift. Take its constant in a horizontal
basis at \(c-1\) and transport it to \(c\) with precision
\(Xq^C\). This gives \(y_1\) in the last Frobenius span,
with \((\gamma-1)y_1\in XN_c(T_1)\) and \(q^C y_1\in N_c(T_1)\).
Use the truncation order
\[
 R_0=\lceil C/(p-1)\rceil
\]
and
take \(\sigma_1=R_0\). In positive degree \(r<R_0\) the
denominator costs \(r+v_p(r!)\). The transported monomial is
a unit times \(X^r q^{(p-1)r}\). Keep one \(X\) and expand
the others as \((q-p)^{r-1}\).

Modulo \(Xq^C\) the surviving
terms contain at least \(\max(0,pr-C)\) powers of \(p\):
a term of extra degree \(j\) must satisfy \(j+(p-1)r<C\).
The loss is therefore at most
\[
 pr/(p-1)-\max(0,pr-C)\le C/(p-1);
\]
the maximum occurs at the crossing \(pr=C\). The transported
tail is zero to this precision. The last-level jet of \(y_1\)
is horizontal in the last Frobenius span: its difference is
zero modulo \(Xq^C\) there, by the truncation and the integral
correction. Over the coefficient field use the local logarithmic
parameter at the primitive level. On a horizontal column the
action on degree \(k\) has eigenvalue \(\chi(\gamma)^k\), so
an invariant jet has only its constant coefficient. This gives
the jet of \(F(y_1)\) from the transported constant of \(n_1\)
with multiplier \(p^{\sigma_1}\).

To compare with \(f_2\Phi\lambda\), use the common horizontal
column at the last level. At this level the transported multiplier
\(\phi(s)=P_2^y/P_1^v\) has constant term
\[
 (P_2^y/P_1^v)(0)=p^{y-v},
\]
since \(P_j(0)=p\) for \(j\ge1\) in the \(q\)-coordinate.
Thus the two scalar constants to be matched, after multiplying
by \(p^{k_0+v}\), are \(x_0p^v\) and
\(p^{k_0+\sigma_1+E+y}\). In the second expression the
common horizontal normalization cancels; its \(p^\sigma\) is
present in both comparisons.

Choose \(t_3,B'\in O\) with
\[
 t_3 x_0p^v=B' p^{k_0+\sigma_1+E+y},\qquad v_p(t_3)\le14L .
\]
For \(x_0=0\) use \((1,0)\). Otherwise take \(t_3=1\) if
\(v_p(x_0)+v\ge k_0+\sigma_1+E+y+\sigma\). In the other
case choose equality with \(B'=p^\sigma\). Equation
\eqref{eq:short-last-constant} bounds \(v_p(t_3)\) by
\[
 \sigma+E+\sigma_1-C/p+C/p^2+2+\gamma_h\le14L .
\]
In fact \(\sigma_1\le C/(p-1)+1\), \(C\le h\), and
\(1/(p(p-1))+1/p^2\le 2.1/p^2\).

For the marked map \(F\), the required Frobenius-span congruence is
\[
 t_3 f_2\Phi\lambda-B' F(y_1)\in q^C\langle\Phi N_{c-1}(T_2)\rangle.
\]

For the first term the numerator at level two supplies the
height loss from the first negative letter; equivalently
\(f_2\Phi\lambda\) is an integral multiple of
\(\Phi(s\lambda)\). For the divisibility assertion compare
the two horizontal jets after inverting \(p\) and the higher
level factors. Modulo \(q^C\) the first has constant multiplier
\(x_0/p^{k_0}\) on \(\Phi\lambda\), by the last congruence;
the horizontal jet of the second has constant multiplier
\(p^{\sigma_1+E+y-v}\). Higher level factors are units at
\(q=0\) over the fraction field, and the bounded integral
Frobenius span is free.

For the free bounded integral Frobenius span \(P\), monic division yields
\[
 P\cap q^C P[1/p]=q^C P.
\]
Thus the congruence of rational horizontal jets is integral. The negative terminal
height puts this difference in \(N_c(T_2)\). Lift \(p^R\)
times the difference there and add \(p^R B' y_1\).
The result is admissible and maps to \(p^Rt_3 f_2\Phi\lambda\);
its denominators in \(N_c(T_1)\) are cleared by \(q^C\).
Since \(\Phi\lambda\) represents \(p^\sigma d\), this
lifts \(t_3f_2\) at cost \(R+\sigma\). The summand \(t_3x_2\)
admits the constant lift at the same cost. The scalar clearing cost is
\[
 \lceil v_p(t_3t_2)\rceil=\lceil v_p(t_3)+v_p(t_2)\rceil.
\]
Adding the lifting costs gives the bound \(22L+1+\eta B\).
\end{proof}


\begin{lemma}[Three-letter lifting with positive outer signs]\label{lem:three-letter-positive-lift}
For the same \(L=S_g\) and short-band indicator \(B\) as in
Lemma~\ref{lem:three-letter-negative-lift}, and for chronological signs \(+,-,+\), the adjacent-level correction can be chosen with exponent \(e_3\) satisfying
\[
 e_3\le34L+1+\eta B.
\] With the other short-arc branches, the
single exponent \(E=\lfloor40S_g+\eta B\rfloor\) works on
every eligible column of the chosen horizontal frames,
in powers of \(\ell\) in the original DVR.
\end{lemma}

\begin{proof}
\proofstep{Use the middle congruence with the two positive neighbor denominators.}
Now consider \(+,-,+\). The denominator of \(f\) is
\(D=q^zP_2^u\) and its middle congruence is modulo
\(P_1^{v+S}\). Write \(C=v+S\), \(a_0=\lceil z/p\rceil\).
We first show that the middle rational constant \(x/p^k\)
has denominator at most \(p^s\), with
\begin{equation}\label{eq:three-denominator}
 s\le\max\{1,u-C/p+16L/5\}+1 .
\end{equation}
Put \(Q=u-(C-a_0+1)/p+2+\log_p u\), and use
\(s=1+\lfloor\max(0,Qp/(p-1))\rfloor\).

If \(D\in(p^s,P_1^C)\), multiply by \(q^{pa_0-z}\).
Expanding \(P_1=p\varepsilon+Uq^p\) gives
\(P_1^{a_0+s-1}\in(q^{pa_0},p^s)\).
Hence \(P_1^{a_0+s-1}P_2^u\in(p^s,P_1^C)\).
Cancelling this power of \(P_1\) modulo \(p^s\), and descending
by \(\phi\), gives
\[
 P_1^u\in(p^s,q^{C-a_0-s+1}).
\]

Cancellation holds by reduction to a distinguished polynomial,
and descent holds since the ring is free over its Frobenius
image (Weierstrass division by \(\phi(q)-Y\)).
The content inequality from the preceding paragraph would force
\[
 s\le\max(0,u-(C-a_0-s+1)/p+2+\log_p u),
\]
contrary to the choice. Here \(C-a_0-s+1>0\):
\(a_0\le z/p+1\), \(20u\le b_0\) and the bound on \(s\)
in \eqref{eq:three-denominator} gives \(a_0+s<C+1\) using
\(160L\le b_0\).

To verify that bound directly, use
\((a_0-1)/p\le h/p^2\), \(u\le w\) and \(\log_p u\le\gamma_h\).
Expanding the factor \(p/(p-1)=1+1/(p-1)\) charges at most
\((w+h/p^2+2+\gamma_h)/(p-1)\) in addition to
\(u-C/p+h/p^2+2+\gamma_h\), which fits the displayed allowance.
This proves the bound on the content of \(D\); the middle
congruence then gives
\[
 p^sx/p^k\in O.
\]


Put \(A=A_{c-2}\). If \(u<A+\eta+2\sigma+2\), separate the
constant at cost \(s\) and lift both summands at \(R+\sigma\).
Indeed the new remainder has numerator divisible by \(P_1^C\).
We check that the size-dependent loss in
\eqref{eq:three-denominator} fits the band allowance.
By \eqref{eq:short-A},
\[
 u-C/p<\eta+\eta/p+\tfrac52 L+2\sigma+2.
\]
Inside the band this gives \(s\le\eta+14L+1\). Outside it,
if \(pu\le b_0\) the excess \(u-C/p\) is nonpositive.
If \(pu>b_0\), absence of the band gives
\(u-b_0/p>2\eta\), while the last upper bound and
\((C-b_0)/p\le(\eta+S)/p\) imply \(\eta<12L\).
These bounds give
\[
 s\le 28L+1+\eta B
\]
on every arc.

If \(u\ge A+\eta+2\sigma+2\), construct a filtered constant
at the initial positive letter. The one-letter calculation with
the coordinate \(q^{-u}d\) gives scalar comparison
\(p^{2\sigma}(q-p)^A q^{G-u}\) modulo \(q^{G-\eta}\).
The degree \(G-u+A\) is in the comparison range. Its leading coefficient is
\[
 [q^{G-u+A}]\,p^{2\sigma}(q-p)^A q^{G-u}=p^{2\sigma}.
\] Equivalently, some coefficient
in that range divides \(p^{2\sigma+1}\): expanding
\(q^{A+2\sigma}\) in \((q-p)+p\) shows that divisibility of all
the low coefficients by \(p^{2\sigma+1}\) is impossible.
All these source coefficients lie in \(\mathcal F^u\). Thus
after transport a source direction at positive level two maps
to \(p^{2\sigma+1}d\). Only its filtration at level two has
been prescribed. We choose the part of \(f\) to which it applies.

Let
\[
 j=\max(0,\lfloor u-C/p-\eta/p^2\rfloor).
\]
There is a
decomposition
\[
 P_2^u=p^j U_0+P_1^C q^\eta U_1,\qquad U_0,U_1\in\mathcal R .
\]
In the expansion of \((p\varepsilon_1+U'_1P_1^p)^u\), the
term with \(P_1^{pr}\) is paid by its scalar coefficient if
\(r\le C/p+\eta/p^2\). Otherwise factor off \(P_1^C\) and
expand the remaining \(pr-C\) powers of \(P_1\).
Any resulting term with \(q\)-degree below \(\eta\) has
scalar valuation at least
\(u-r+(pr-C)-\eta/p\ge u-C/p-\eta/p^2\).
This proves the decomposition. Put \(e_4=(s-j)_+\) and
\(x'=p^s x/p^k\).

Take
\[
 f_a= p^{e_4+j-s}x' U_0/P_2^u.
\]
Then \(p^sf_a-p^{e_4}x'\) belongs to
\(P_1^Cq^\eta P_2^{-u}\mathcal R\).
The second rule lifts \(f_a\) at cost \(3\sigma+1\), using
the source filtration at level two and the constant congruence
at the omitted positive level zero. The remainder
\(p^{e_4}f-f_a\) has numerator divisible by \(P_1^C\) over
\(D\), by the middle congruence and saturation. It costs
\(R+\sigma\). Since \(j\ge u-C/p-\eta/p^2-1\),
\eqref{eq:three-denominator} gives
\[
 e_4\le 16L/5+\eta/p^2+2
\]
when the second term in the maximum is used; otherwise
\(e_4\le2\). In particular \(R+3\sigma+1+e_4\le18L+1\).
The two initial-positive cases therefore both fit
\(34L+1+\eta B\). \label{cor:inert-short-positive-1}

All the bounds for the individual scalar lifts are at most
\(40S_g+\eta B\), since \(S_g\ge1\).
They hold on every eligible column of every chosen horizontal
frame; consequently the integer
\(E=\lfloor40S_g+\eta B\rfloor\) can be used simultaneously.
The exponents are in powers of \(\ell\) in the original DVR,
including when its uniformizer has fractional normalized valuation.
\end{proof}


\begin{lemma}[A fixed exceptional span for a short arc]\label{lem:short-arc-exceptional-span}
For the short arc of length \(a\) in the cut comparison,
one may choose an exceptional set \(Z_c\) of target classes, independent of the input, with
\[
 \#Z_c\le(a+1)ma.
\] Outside its span the scaled lift has
exponent \(\sigma+E+\kappa_c\). The signed count is
\(r_c=a\) because each letter has an off-central window.
\end{lemma}

\begin{proof}
\label{par:short-exceptions}
On the marked target over the coefficient field let \(U\) be the
intersection of the transported \(\mathcal F_j^{u_j}\) over
positive levels. Let \(V_-\) be the sum of
\(\mathcal F_j^{-\eta}\) over negative levels. By the bound on
central multiplicities,
\[
 \operatorname{codim}U+\dim V_-\le ma .
\]
In the initial integral basis, the constants form the lattice
\(D_O\) with the transport of scalars to the terminal crystalline
space understood. Split the saturated sublattice \(D_O\cap U\)
off with complement \(K\), and project \(V_-\) along \(K\)
to \(U\). The intersection of the projected space with the
lattice of \(U\) is again saturated. Use adapted bases of these
two splittings.

Let \(e_j\) denote the adapted basis and let \(J\) index the complement \(K\) and this last intersection. Then
\[
 |J|\le ma,\qquad V_-\subseteq\operatorname{span}_{K_O}\{e_j:j\in J\}.
\]
Every other column belongs to \(U\).
The two changes of basis are integral and invertible, so they
give a Wach frame at \(c'\). Take the unique horizontal columns
with these initial values.

Apply the cut construction of Section~\ref{sec:inert-long-orbit}
to a crystalline extension representing a finite target class.
Its cut vector is the difference in that extension of a vector
in the incoming Frobenius span and a terminal Wach vector.
At each cyclotomic level the incoming span is regular; a pole
of order \(k>0\) therefore has coefficient in the extension's
\(\mathcal F^k\), and its image lies in the kernel of projection
to the trivial representation. Strictness identifies this
intersection with \(\mathcal F^k\) of the target.

For a polar term \(a_{-k}t^{-k}\) with \(k>0\), strictness therefore gives the two respective tests
\[
 a_{-k}\in V_-\quad(\text{negative letter}),\qquad k\le u_j+S\quad(\text{positive letter}).
\]
The statement for an entire polar part is obtained
coefficient by coefficient in the Laurent expansion, followed
by Chinese remaindering on the finite jets. At negative levels
all negative degrees thus disappear on indices outside \(J\).
Likewise \((\gamma-1)w\in XN_c\) imposes the Wach filtration
on the divided differences. Their projection outside \(J\)
vanishes at a negative level in degrees below \(u_j\):
for lower degrees the required filtration is contained in
\(\mathcal F_j^{-\eta}\), or equals it by the gap.

Write \(\mathcal J_c\) for the union of signed degree intervals
used in \(\kappa_c\).
Apply to these horizontal coordinates the integral operator
\[
 K(\gamma)=\left(\sum_{j=0}^{p^a-1}\gamma^j\right)
     \prod_{k\in\mathcal J_c\setminus\{0\}}
            (\gamma^{p^a}-\chi(\gamma)^{kp^a}) ,
\]
On a Laurent coefficient of
degree \(k\) at level \(t<a\), the principal action raised to
\(p^a\) is the scalar \(\chi(\gamma)^{kp^a}\). The sum
projects degree zero to a common value on the primitive roots
after tame invariants. Outside \(J\) this removes the excess
positive-level poles and the negative-level degrees from \(u_j\)
through \(u_j+S-1\).

At a negative level the divided difference condition on a nonconstant coefficient \(a_k\) below \(u_j\) reads
\[
 (\chi(\gamma)^{kp^a}-1)a_k=0,\qquad k\ne0.
\]
The scalar is nonzero, hence \(a_k=0\).
Thus the transformed scalar coefficients outside \(J\) satisfy
\eqref{eq:short-scalar}. The remaining constant at each negative
level is in the coefficient field, so a scalar power clears it.
The operators preserve the bounded integral denominator
\(\Lambda\); the improved pole orders descend to \(D\) by
division by the distinguished factors. Outside \(J\), the only poles allowed by \(D\) are positive with directions in \(U\); the negative zeros have the required orders. Write \(f=f_J+f_{J^c}\) for the support decomposition and \(\Delta\) for divided difference. The Wach test is linear, so
\[
 \Delta f_J=\Delta f-\Delta f_{J^c}
\]
passes as well.

For the latter part let \(M_J\) consist of the tuples with
denominator \(\Lambda\) supported on \(J\); let \(I_J\) be
the submodule whose divided differences pass the Wach test.
Let \(P_J\) consist of \(X\) times the tuples passing the Wach test
with bounded denominator, in \(M_J\). For numerators supported on \(J\),
a change by \(X\Lambda^2h\) induces the following change of tuple:
\[
 \frac{f+X\Lambda^2 h}{\Lambda}-\frac{f}{\Lambda}
       =X\Lambda h\in P_J .
\]
Weierstrass division reduces each numerator modulo \(X\Lambda^2\);
hence \(M_J/P_J\) is finite over \(O\).

The quotient \(M_J/P_J\) is also torsion-free.
Suppose \(0\ne c\in O\) and \(c(f/\Lambda)\in P_J\). Monic division by
\(X\) gives \(cf(0)=0\), hence \(f(0)=0\). The tuple
\(f/(X\Lambda)\) now has bounded denominator. Its rational filtration
tests hold since \(c f/(X\Lambda)\) satisfies them and the filtration
spaces are vector spaces over the coefficient field. Thus
\[
 f/\Lambda\in P_J.
\]


Let \(E_J\) be the image of \(I_J\) in \(M_J/P_J\).
After inverting \(p\), Chinese remainders for \(X\) and the
distinct cyclotomic factors of \(\Lambda\) split the numerator
jets at \(0\) and at the \(a\) levels. At \(0\) the quotient by
\(X\) leaves at most one entry per coordinate in \(J\).
At a cyclotomic level \(X\) is invertible, so the quotient is
by the filtered Laurent lattice.

In coefficient degree \(k\),
the principal power fixing the roots acts on that quotient by
\(\chi(\gamma)^{k p^a}\). A nonzero fixed class therefore needs
\(k=0\), since \(\chi(\gamma)\) has infinite order. In degree
zero the tame and principal root actions are transitive. Invariants
have at most one copy of each supported coordinate, and passing
to the quotient cannot increase this invariant dimension by
the characteristic-zero averaging projection. Thus
\[
    \dim_{\operatorname{Frac}(O)}(E_J[1/p])
          \le |J|+\sum_{t=0}^{a-1}|J|=(a+1)|J| .
\]
Since \(M_J/P_J\) is finite torsion-free over \(O\), this is
also a bound on the number of generators of \(E_J\).

Take an \(O\)-basis of the image and choose its representatives
in \(I_J\) once for the cut. Realize \(p^\sigma\) times each
by the horizontal comparison and take its cut Herr class. The
resulting set \(Z_c\) consists of finite classes, since a
product of cyclotomic factors carries the representative into
the Wach lattice. Using one representative per basis vector of \(E_J\) gives the uniform bound
\[
 |Z_c|\le\operatorname{rank}_O E_J\le (a+1)ma.
\]

\proofstep{Reduce every class against the same chosen representatives.}
For any input the exceptional coordinates are an integral
combination of these representatives modulo \(P_J\).
Terms of \(P_J\) represent elements of \(XN_c\) after the
same comparison and give zero under the cut class map.
The other coordinates lift at the common cost \(E\). To see the
exponent accounting, the inverse horizontal comparison reads
coordinates of \(p^\sigma w\). Realizing the split coordinates
requires another \(p^\sigma\), while on the outside part the
vector lifting relation already represents \(p^E\) times the
coordinates.

Since \(E\ge\sigma\), pad the realized split parts
by \(p^{E-\sigma}\). Their sum now matches
\[
 p^{\sigma+E} K(\gamma)w
\]
modulo \(XN_c\). Naturality and the vanishing on
\(XN_c\) follow from Frobenius contraction in the Herr complex,
as in the cut construction. Write \(K(T)-K(1)=(T-1)Q(T)\). Stability and the cut-vector condition give
\[
 K(\gamma)w-K(1)w=Q(\gamma)(\gamma-1)w\in XN_c.
\]
Thus the two vectors give the same cut class. In the scalar reading
\[
        p^{\sigma+E}K(1)w
             = u_c p^{\sigma+E+\kappa_c}w,\qquad u_c\in O^\times
\]
since \(v_p(K(1))=\kappa_c\). This gives the exponent
\(\sigma+E+\kappa_c\) with the same set \(Z_c\) for every
input class. For the present signed arcs every letter has a
positive or negative off-central window, so the signed level count
\(r_c=a\). \label{cor:inert-short-exceptions-1}
\end{proof}



\paragraph{Adding the short-arc colengths.}
\begin{lemma}[The common colength from the short cuts]
\label{lem:ordinary-inert-short-cuts}
In either ordinary inert range with a position having smaller jumps, use its
marked arcs. The normalized finite-class colength is bounded by
\eqref{eq:nonthick-colength}; in the range requiring exceptional representatives, there is one set \(Z\) for all the comparisons such that
\[
 \#Z\le4nm.
\]
\end{lemma}

\begin{proof}
Apply Lemma~\ref{lem:short-arc-exceptional-span} with \(p=\ell\) and \(a=a_c\).
On the finite jets at levels \(t<a_c\), the action of \(\gamma^{\ell^{a_c}}\)
on a Laurent coefficient \(c_k\) of degree \(k\) is multiplication by
\(\chi(\gamma)^{k\ell^{a_c}}\). For each specified nonzero degree
the corresponding factor of \(K(\gamma)\) therefore satisfies
\[
 \bigl(\gamma^{\ell^{a_c}}-\chi(\gamma)^{k\ell^{a_c}}\bigr)c_k=0 .
\]

To estimate \(\kappa_c\), count these degrees with their signs. For a positive
jump window these degrees lie in \([-s_i-S,-s_i)\) and for a
negative window in \([s_i,s_i+S)\). Thus there are at most
\(a_c S\) degrees, each nonzero with absolute value at most \(h\).
In any one of these intervals of \(S\) integers the count of
multiples of \(\ell^j\) is at most \(S/\ell^j+1\).
For such an interval \(I\) excluding zero the resulting valuation bound is
\[
 \sum_{k\in I}v_\ell(k)=\sum_{j=1}^{\lfloor\log_\ell h\rfloor}\#\{k\in I:\ell^j\mid k\}\le S/(\ell-1)+\log_\ell h.
\]

The principal trace contributes at most \(a_c\), and the
nonzero degrees \(k\) contribute at most \(a_c+1+v_\ell(k)\) each,
by the principal-unit valuation used in
Section~\ref{sec:inert-long-orbit}. Hence
\[
 \kappa_c\le a_c+(a_c+1)a_c S
                +a_c(S/(\ell-1)+\log_\ell h)
            \le13 S+3(1+\gamma_h).
\]
Given the cut decomposition and the arc lift, pad the powers to
their common maximum and sum the lifted cut classes.
The set \(Z\) of exceptional classes
is contained in the union of the sets from the cuts, so
\[
 |Z|\le\sum_c m r_c(a_c+1)\le4m\sum_c a_c=4nm .
\]
The cost of the scaling on the original class is at most
\(55S_g+13S+\eta B\), since \(1+\gamma_h\le S_g\).
Thus the length argument of Section~\ref{sec:inert-thick-word}
applies to the same integral finite group.

In the larger range
one uses the full lift and takes \(Z\) empty. In both ranges the
normalized colength is at most
\begin{equation}\label{eq:nonthick-colength}
 320(n+1)r\bigl(h/\ell^2+(w+\eta)/\ell+
                      S+1+\gamma_h+\eta B\bigr)+320nm e' .
\end{equation}
Indeed \(S/\ell\le S\), all the terms in the parentheses
are nonnegative, and \(55+13\le320\). Lemma~\ref{lem:short-arc-exceptional-span}
supplies the fixed representatives for this estimate. The coarse exponent
\(e'\) is for the actual marked surjection and is used only on the
image of the span of the exceptional finite classes.
\end{proof}

\begin{proposition}[One normalized entry for all ordinary inert rows]\label{prop:ordinary-inert-place-entry}
In the ordinary inert range, with the marked word data and
\(U_H,U_F,U_M\) defined in Section~\ref{sec:cut-length-rows},
there is one rational place entry \(c_w\) for every good character at the place. Its range is
\[
 0\le c_w-(368U_H+2082U_F+22365U_M)\le U_F.
\]
Here the indicator entering \(U_F\) is the common place indicator,
and every normalized length uses the same coefficient valuation.
\end{proposition}

\begin{proof}
The common unramified coarse finite-class lifting theorem uses the
negative and positive bounds \(M_i,H_i^+\) at each embedding.
Put \(M_*=\max_i M_i\), \(H'=\max_i(M_i+H_i^+)\)
and \(a_*=\lfloor M_*/(\ell-1)\rfloor+1\).
If \(J\) is the set of nonzero degrees formed by adding
\(1,\ldots,M_*\) to jumps and also including \(1,\ldots,M_*\),
the exponent for this theorem is
\[
  2\left\lfloor\frac{H'\ell}{(\ell-1)^2}+3+
                        \log_\ell(H'+1)\right\rfloor
       + a_*+v_\ell(a_*!) +\sum_{j\in J}v_\ell(j).
\]
The theorem supplies a finite lift for every target class after
multiplication by this power of \(\ell\).

In the general boxes \(J\subset[-H',2H']\setminus\{0\}\).
Legendre's formula bounds the valuation sum by \(3H'/(\ell-1)\).
Also
\(a_*+v_\ell(a_*!)\le(\ell/(\ell-1))(H'/(\ell-1)+1)\).
Using \(\ell/(\ell-1)\le11/10\) gives
\[
 e'\le23\bigl(2+H'/\ell+\log_\ell(H'+1)\bigr).
\]
For the large-\(S_g\) boxes above, substitute \(J\subset[-h,2h]\),
\(M_*\le h\), and \(H'\le h+\eta\le21h/20\) into the same
expression. Adding the coefficients of \(h/\ell\) from the three terms gives
\[
 2(21/20)(11/10)^2+(11/10)^2+33/10<7.1.
\]

The constants other than the logarithm total at most \(7.1\),
and \(H'+1\le2h\) since \(H'\) is integral and \(h\ge1\).
This proves the exponent estimate used in that range.

Use the global \(H,w_P,S_c,\eta_3,z\) and the quantities
\(U_H,U_F,U_M,\mathcal P_\ell\) of
Section~\ref{sec:cut-length-rows}. The local transport gives
\(S\le S_c\), \(\eta\le\eta_3\), \(w\le w_P\), and \(h\le H\)
in both branches. In the case in which every position has larger jumps this last inequality
means \(\max_i b_i+S\le H\). The polynomial size bounds give
\(1\le H\le2e^z\), \(z\ge1\) and \(\eta_3\le H\).
When some jumps are in the smaller range, \(H'\le2H\). In the
case in which all jumps are in the larger range
\[
 H'\le H_{\rm loc}\le H+\eta_3\le2H.
\]

As \(H'+1\le3H\le\ell H\) for the logarithm's argument, it follows
that \(e'\le69\mathcal P_\ell\).

For \eqref{eq:nonthick-colength}, \(\gamma_h\le3+z\) follows from
\(2h\le4e^z\) and \(\log\ell\ge1\).
Consequently \(1+\gamma_h\le5z\). The row length is at most
\(320U_H+1600U_F+22080U_M\).

For Proposition~\ref{prop:thick-inert-colength},
\(L_0\le8+z\le9z\), since
\(2H_{\rm loc}\le8e^z\). Further,
\[
 32H_{\rm loc}/\ell^2\le32H/\ell^2+3\eta_3/\ell,\qquad
 10\eta/(\ell-1)\le11\eta_3/\ell .
\]
The exponent \(e\) is therefore at most
\(32H/\ell^2+225F_{\ell,\lambda}\).
Also \(H_{\rm loc}+1\le3H\), which gives
\(c_D\le4\mathcal P_\ell\).
Substituting in \eqref{eq:thick-chain-mass} bounds the colength
by \(32U_H+227U_F+288U_M\).
Both branches thus have common majorant
\[
                       368U_H+2082U_F+22365U_M .
\]
The place entry is a rational number between this value and its
sum with \(U_F\). Here \(U_F>0\) by \(z\ge1\) and \(r>0\).
All comparisons use the same place indicator and coefficient
normalization. Hence the entry bounds the row for every
good character at once.
\end{proof}

